% SPDX-License-Identifier: CC-BY-4.0
% Copyright (c) 2026 Jonathan Yantis. Research draft; not submitted.
\documentclass[11pt,letterpaper]{article}
\usepackage[margin=0.85in]{geometry}
\usepackage[T1]{fontenc}
\usepackage{mathptmx,amsmath,amssymb,booktabs,array,ragged2e,microtype,xcolor,xurl,hyperref,fancyhdr}
\hypersetup{colorlinks=true,urlcolor=blue!45!black,citecolor=blue!45!black,linkcolor=black,pdfauthor={Jonathan Yantis},pdftitle={PToB Viability: Rooted BLE Timing Meshes for Heterogeneous Edge Networks}}
\setlength{\parskip}{3pt}\setlength{\headheight}{14pt}\setlength{\emergencystretch}{3em}
\widowpenalty=10000\clubpenalty=10000
\pagestyle{fancy}\fancyhf{}\fancyhead[L]{\small PToB viability research preprint}\fancyhead[R]{\small 30 October 2026}\fancyfoot[C]{\thepage}
\renewcommand{\topfraction}{0.9}\renewcommand{\textfraction}{0.1}\renewcommand{\floatpagefraction}{0.8}
\newcolumntype{L}[1]{>{\RaggedRight\arraybackslash}p{#1}}
\newcommand{\us}{\,\mu\mathrm{s}}\newcommand{\ns}{\,\mathrm{ns}}
\newcommand{\repo}[2]{\href{https://github.com/yantisj2/PToB/blob/a3c96db35b1817de534abe4450cfd032ceca77e8/#1}{#2}}
\title{\textbf{PToB Viability: Rooted BLE Timing Meshes\\for Heterogeneous Edge Networks}\\[5pt]\large Causal Time Delivery Across Radios, Relays, and Hosts\\Architecture, Evidence, and Operating Limits}
\author{Jonathan Yantis\\\small Independent engineering research}
\date{Research draft for author review --- 30 October 2026}
\begin{document}
\maketitle
\begin{abstract}
Precision Time over BLE (PToB) turns periodic radio observations into a clock service spanning heterogeneous microcontrollers, regenerating relays and host applications. Its design keeps source permission, causal clock-map validity and child admission independent: receiving a fresh packet or predicting a stable clock cannot renew expired authority. Relays transfer identified physical events into a root coordinate, while capability metadata distinguishes actual capture from estimated transmission. We contribute this service composition, its cross-platform implementation and a reproducible synthesis of distinct experiments through 30 October 2026. A W12/T096 relay study measured 60.328 ns relative PPS RMS on 522 of 600 logged receiver records and 61.922 ns on 451 of 600 downstream records. A separate N30 experiment retained a 1591 ns maximum despite 166.590 ns RMS. Frozen BLE-trained W12 clocks supported a best 64 s rolling sampled sub-microsecond P99 duration; a distinct fixed cohort passed 75 s and failed at 76 s, while later windows shortened. USB-to-macOS tests extend relative time into an application clock with separate prediction, agreement and availability populations. These interface results support shallow managed timing cells rather than a qualified composed deployment. Timing, age, airtime, availability and energy models identify the scaling constraints. Public selected data preserve failures and permit offline numerical review. Calibrated UTC/electrical-output accuracy, physical failover and sleep energy remain open. The immediate opportunity is useful event time across diverse endpoints with explicit provenance and refusal, with PTP ingress and application adapters as proposed extensions.
\end{abstract}
\noindent\textbf{Keywords:} clock synchronization; Bluetooth Low Energy; boundary clocks; rooted mesh; hot standby; sensor networks; PTP; edge computing.
\par\noindent{\footnotesize Copyright \copyright\ 2026 Jonathan Yantis. Manuscript and figures: \href{https://creativecommons.org/licenses/by/4.0/}{CC BY 4.0}. Research draft; not submitted or peer reviewed.}

\section{Introduction}
Edge devices need a shared interpretation of event time without a timing cable, GNSS antenna or complete Ethernet timing interface at every endpoint. Mobile sensors, radio-to-host instruments and relays around coverage obstacles motivate distributing an identified clock coordinate over inexpensive local radio infrastructure.

The engineering question is whether useful time survives those boundaries with explicit validity and uncertainty. A plausible display after source loss is insufficient. Relaying an old timestamp without accounting for residence, oscillator aging and path changes can extend apparent coverage while degrading temporal meaning.

PToB addresses this question with a rooted distribution architecture. A small source-control plane selects and authorizes a serving grandmaster (GM) and, when prerequisites hold, an actually synchronized hot standby (HOT). A periodic timing plane supplies physical-event observations and subsequent corrections. Relay clocks (RCs) maintain their own causal clock maps and generate fresh downstream events. Clients consume those maps at different capability levels, including a macOS application clock reached through a USB timing bridge. The network can have redundant radio reachability while each active timing path remains rooted and acyclic.

The distinctive solution is a common clock-service discipline across boundaries that are often evaluated separately. A relay has to preserve the root's permission while regenerating its own event; a host has to preserve the source map while crossing USB and a general-purpose OS; a held clock has to distinguish prediction skill from permission to serve. PToB makes these obligations explicit rather than reducing each boundary to a fitted offset.

We contribute three connected results: (i) a concrete BLE service composition with separate origin permission, clock-map validity and admission; (ii) causal event-to-root transfer across heterogeneous regenerating clocks and a host bridge, with explicit capture capabilities; and (iii) public, failure-aware interface evidence, including sampled prefix-maximum holdover. Conditional scaling models and a proposed PTP ingress identify how this service could extend. The hypothesis is that shallow managed BLE cells can deliver useful shared event time across diverse endpoints. The experiments test individual interfaces; composed end-to-end qualification remains a separate step.

The evidence cutoff is 30 October 2026. This revision incorporates public selected observations and research completed since the 16 October edition, which remains preserved in Git \cite{legacy}. No new hardware experiment was performed to write this synthesis. Historical raw-radio and fixed-route measurements, later GM7/RC7 implementations, ESP32-S3 captures, USB/chrony tests and holdover cohorts retain separate identities. This edition is an author-review draft prepared toward arXiv submission.

\section{The service being evaluated}
\subsection{Position within wireless time synchronization}
PToB builds on timestamping and skew estimation established by FTSP \cite{ftsp} and BLE synchronization studied by CheepSync, BlueSync and MicroSync \cite{cheep,bluesync,micro}. Periodic-advertisement energy and relay studies motivate the bearer \cite{paenergy,parelay}; PTP boundary clocks, gPTP standby/domain-quality selection, secure-time authorization and recursive trust, and rooted routing supply architectural precedents \cite{ptp,gptpstandby,securetime,rpl}. The contribution is their concrete BLE realization across source control, causal clock delivery, heterogeneous regeneration and a host boundary, with inspectable validity and failure accounting. Different instruments, cohorts and power boundaries preclude a numerical ranking.

The service uses BLE periodic advertising \cite{ble}; it is a project-specific protocol rather than IEEE 1588 or Bluetooth Mesh. PTP and the SMPTE broadcast profile motivate proposed ingress and media adapters \cite{ptp,smpte}. Qualification of those adapters is separate from the BLE timing results.

\subsection{A multidimensional viability test}
We use five criteria. \emph{Temporal validity} asks whether delivered time refers to an identifiable source, physical event, map generation, and permitted interval. \emph{Precision and calibration} ask whether useful relative phase is achieved and which systematic uncertainties remain. \emph{Availability} includes startup, expiry, reference loss, route changes, and recovery. \emph{Scalability} includes control state, radio occupancy, receiver scheduling, and path depth. \emph{Economics} includes hardware, installation, energy, maintenance, and qualification cost.

A platform can be viable for timestamping intermittent sensor events while unsuitable for continuous actuator synchronization. Application requirements must therefore specify the observation point, maximum error or statistical objective, required availability, outage behavior, and trusted source. An advertised ``sub-microsecond'' value without these attributes is insufficient.

\section{Protocol architecture and invariants}
\subsection{A clock read is a service decision}
PToB's common abstraction is a causal map plus the conditions under which its result may be consumed. A read binds a source/path scope, map generation, local and inherited age, original deadline, capture capability and known or unknown uncertainty. These describe different properties: a fine counter improves representation, an accurate forecast describes prediction, and authenticated source permission governs service. Heterogeneous adapters preserve those meanings even when their capture mechanisms differ.

For an admitted consumer, the intended publication rule can be written
\begin{equation}
 V(t)=P_{\mathrm{source}}(t)\land M_g(t)\land A_{\mathrm{path}}(t),
\end{equation}
where $P_{\mathrm{source}}$ is current origin permission, $M_g$ is the causal map's validity and $A_{\mathrm{path}}$ is admission in the selected scope. This is a service-validity predicate, not an accuracy bound. An unexpired map with unknown physical uncertainty still has unknown accuracy; a numerically stable map with expired permission cannot authorize downstream output. A local HOLD estimate can remain useful under its declared policy while network serving is withdrawn.

Consider a relay and USB client during a source change. Renewal in the same scope may preserve still-valid captures and maps, but does not renew origin age. A changed root/boot/path fences incompatible history and requires fresh admission where applicable. The relay's new event and the host's later query must use maps already available at their evaluation instants. Neither a queued correction nor a host daemon may silently revive the old scope. This cross-boundary invariant is the organizing principle of the design; the evidence below tests its interfaces separately.

\subsection{Rooted timing over a potentially redundant mesh}
The logical timing graph is a tree or directed acyclic rooted subgraph selected from the available radio graph. A receiver disciplines to one selected parent in a specific source scope. Alternative parents provide recovery choices rather than clocks to average indiscriminately. The root represents independent source authority; an RC extends its time without acquiring GM entitlement. A leaf can remain receive-only during steady timing service, although admission or renewal may require occasional interaction.

Figure~\ref{fig:arch} separates demonstrated interfaces from the proposed PTP ingress. Neither a desktop application nor an unqualified relay can silently promote an uncertain local estimate to an independent time source.
\begin{figure}[htbp]\centering
\setlength{\unitlength}{1pt}
\begin{picture}(460,178)
\put(0,133){\framebox(120,34){\shortstack{PTP / facility backbone\\\scriptsize proposed ingress}}}
\put(150,133){\framebox(130,34){\shortstack{Authorized source / GM\\\scriptsize GNSS demonstrated}}}
\put(315,133){\framebox(140,34){\shortstack{Eligible HOT standby\\\scriptsize bounded field evidence}}}
\put(120,150){\vector(1,0){30}}\put(280,150){\vector(1,0){35}}
\put(215,133){\vector(0,-1){31}}
\put(150,68){\framebox(130,34){\shortstack{Regenerating relay(s)\\\scriptsize one selected timing parent}}}
\put(65,68){\framebox(64,34){\shortstack{Leaf\\\scriptsize sensor}}}
\put(150,85){\vector(-1,0){21}}
\put(315,68){\framebox(140,34){\shortstack{USB timing receiver\\\scriptsize source-scoped bridge}}}
\put(280,85){\vector(1,0){35}}
\put(385,68){\vector(0,-1){29}}
\put(315,5){\framebox(140,34){\shortstack{macOS application clock\\\scriptsize private chrony, no OS steering}}}
\put(0,13){\makebox(290,30)[l]{\shortstack[l]{\small Timing paths preserve origin, age and permission.\\\small Arrows show logical flow, not a single tested topology.}}}
\end{picture}
\caption{System architecture. The relay, HOT, and host results come from separate experiments; the complete drawing is not a measured combined deployment.}\label{fig:arch}
\end{figure}

\subsection{Physical events and causal clock maps}
A timing exchange proceeds as follows \cite{current,v14}: the receiver saves an original's physical RX capture with sender/stream incarnation, event identity and channel/capture scope. A later matching correction supplies that event's TX coordinate. After authentication, association and declared delay treatment, processing can activate a new immutable map. Stale, conflicting or ambiguous matches are refused. The relay then maps its own next transmit event through the generation active at that TX and publishes a corresponding follow-up under the inherited source deadline. Actual and estimated TX backends retain their declared capability.

For local counter $c$ and immutable generation $g$, a generic map is
\begin{equation}
\widehat T_g(c)=a_g(c-c_g)+b_g.
\end{equation}
The implementation may use fixed-point arithmetic and a slewed actuator. For a live read or prospective witness score, generation $g$ must have been active at that read/event. A delayed event estimate may use later information with its availability time declared; it cannot establish what the earlier running clock delivered. New maps activate only after actual processing completes. Handover preserves the declared continuity policy or withdraws service.

Regeneration permits downstream periodic service between upstream observations while preserving origin age and permission. Its event-to-root mapping replaces forwarding with a guessed residence delay. An estimated-TX relay can provide functional service while its physical timing uncertainty remains unknown.

\subsection{Four distinct lifetimes}
The implementation experience strongly favors separating (1) authenticated control liveness, (2) independent source permission, (3) local timing-map validity, and (4) child service admission. A live peer can temporarily lack a precision map; a long child lease cannot extend short source permission. Freshly receiving old provenance cannot reset its age. Display smoothing cannot extend the API's clock validity.

A published sample therefore carries or binds root identity, physical boot/source epoch, selected path, sender/stream incarnation, map generation, and original deadlines. Changes fence incompatible history and cancel queued delivery where appropriate. A ``HELD'' output is meaningful only within a declared valid hold interval. Missing, unknown, withdrawn, and expired states remain distinct; none becomes valid because the previous RMS was small.

For example, the fixed-route research profile uses local map age strictly below 500 ms. Its admission representation can last up to 360 s, while source permission is much shorter \cite{v14}. These are different clocks in the protocol state machine, not alternative precision guarantees. Inherited source age must continue to accumulate across regenerating relays.

\subsection{Source selection and hot standby}
Source eligibility precedes ranking: good mesh-derived time cannot make a node an independent GM. Current GM7 uses authenticated broadcast discovery and ordinary status observations; GMs follow the selected timing stream without the client JOIN exchange. Ordinary RCs still require admission. The ACTIVE source nominates an independently qualified standby; HOT confirmation requires actual same-scope shadow readiness. Functional readiness and the stricter precision-lock label are separate \cite{current}. Loss of precision lock alone is not evidence that the primary has died.

HOT is an actual operating condition: the selected standby is synchronized and has current takeover permission. A standby label, capable hardware, or an election preference is insufficient. The purpose is to reduce recovery work by preparing a replacement, not to promise zero phase discontinuity or immediate failover. Independent source loss, takeover fencing, source incarnation, and downstream reacquisition still matter.

Disconnected partitions cannot infer global source uniqueness from local radio evidence. Current GM7 retains mesh birth, founder and leader-term lineage; on meeting, the newer ACTIVE mesh yields to the older. This specifies reconciliation and fencing, not measured convergence under arbitrary asymmetric reception. The designated HOT has a conditional three-second first-G1 takeover target, rather than a demonstrated loss-to-output guarantee \cite{current}. Physical stop-before-start, continuity and no-HOT recovery still need instrumented tests. Historical three-GM instability remains part of the evidence; it neither certifies nor automatically diagnoses the successor implementation.

\subsection{Authentication and heterogeneous capability}
The profiles use truncated HMAC-SHA256 with provisioned group credentials and separate authentication domains for control, origin and path records \cite{v14,current}. This provides group authentication; any holder can forge statements in its domain. Boot/epoch, sequence, challenge, path and original deadlines constrain replay after admission, but cannot prove the age of a previously unseen advancing recording at first contact or after replay-state loss. Authenticated one-way timing also retains a fundamental delay ambiguity \cite{delaysecurity}. Revocation, compromise recovery and delay resistance remain security-profile work. Receivers discover meshes locally, keep credentials separate from runtime selection and authenticate admission afresh after a context change. Authenticity establishes neither time accuracy nor resistance to jamming.

Heterogeneity should be explicit in capabilities: actual versus estimated event capture, counter resolution, clock-map availability, calibrated delay, output interfaces, oscillator control, and uncertainty state. A coarse display client, an actual-event receiver, a software-estimated transmitter, and a hardware-timed output device can share source provenance without sharing a precision claim. Radio adapters must preserve supported controller ownership, event identity, and bounded work; moving a timestamp into application code is not automatically equivalent to a hardware capture.

\section{Evidence synthesis and measurement boundaries}\label{sec:evidence}
\subsection{Method}
We reviewed dated manuscripts, selected observations, protocol summaries and indexed provenance. The edition manifest binds the manuscript/PDF and public claim inputs; offline verifiers recompute the selected CSV and frozen-clock populations. Earlier HOT and demand-driven relay results are retained summaries from the original synthesis, whose underlying private captures are not publicly reconstructible. No unlike runs are pooled or interpreted as controlled platform comparisons. RMS includes signed bias; quantiles retain each study's convention (linear interpolation for historical V14, nearest rank for the newer populations). Estimator residual, causal forecast, local PPS witness and electrical-output accuracy remain different metrics. Point samples do not establish continuous coverage.

Table~\ref{tab:evidence} summarizes complementary timing results. A usable result must pair its conditional error distribution with a declared opportunity ledger: eligible, invalid, withheld and missing. For a stated error budget $B$, a conservative observed success fraction is the number of valid scored opportunities with $|e|\leq B$ divided by all declared opportunities; unobserved errors remain unknown. This is neither a per-link packet probability nor an interpolated time-availability estimate.
\begin{table}[htbp]\centering\small
\caption{Separate timing populations; RMS includes bias. PPS denotes relative local witness phase, not measured electrical output. No row establishes calibrated absolute UTC.}\label{tab:evidence}
\begin{tabular}{@{}L{0.22\linewidth}L{0.36\linewidth}L{0.36\linewidth}@{}}\toprule
Boundary / date & Observation & Population and limit\\\midrule
Nordic relay, 12 Oct. & 62.504 ns RMS; 191 ns maximum; 601.576 s & 436/602 nominal PPS opportunities; reference gaps and map expiry retained\\
N30 fine capture, 22 Oct. & 166.590 ns RMS; 1591 ns maximum; 660 s & 657/660 eligible; three invalid, none nominally missing; empirical RX convention\\
W12 reception, 24 Oct. & 60.328 ns RMS; 180 ns maximum; two 300 s windows & 522/600 logged PPS eligible; 78 source-unavailable; frozen empirical calibration\\
T096 via W12, same study & 61.922 ns RMS; 211 ns maximum & 451/600 logged PPS eligible; 109 no-map, 36 expired, four reference-ineligible\\
W12 emitted timestamps, same study & 108.306 ns RMS; 226 ns maximum & 58/120 selected downstream RX comparisons; 61 missing matches, one unbracketed; TX estimated\\
USB-to-host, 17 Oct., separate RC / DC & Causal forecast RMS 190.050 / 267.448 ns; shared-source clock agreement 325.545 / 344.546 ns & 5054 / 5370 forecasts of 5760 USB events each; 145 / 119 clock comparisons; shared bus, separate source chains\\\bottomrule
\end{tabular}
\end{table}

\subsection{Native clock discipline and one-relay timing}
Historical software-clock experiments measured 69.136 and 68.106 ns causal RMS on 1144/1144 and 727/728 quiet-pair originals. A separate three-unit capture measured 1075.679 ns on 354/384 scoreable originals, with no qualified native lock among 59 status observations \cite{osc}. These demonstrate operating-condition sensitivity in the earlier raw-radio profile, not performance of the later periodic controller.

The prospective relay experiment used a source GM, one dedicated RC and a leaf's local PPS witness. That witness was excluded from the delivered estimator and peer acquisition aid. A separate direct-path population fixed a provisional $+10543\ns$ correction per received BLE link; the two-link correction was $+21086\ns$ \cite{v14}.

Over 601.5761865 s, 504 physical PPS records were retrieved. There were 436 reference-valid causal comparisons, with mean $-18.417431\ns$, uncentered RMS $62.503871\ns$, and demeaned jitter $59.728822\ns$. Absolute P95/P99/maximum were 123.25/160.25/191 ns. Sixty-five observed edges lacked valid witness metadata and three reference-valid edges had an expired delivered map. Using 602 conservative nominal boundaries yields 72.425\% usable coverage, including unobserved boundaries. The largest witness gap was 59.017190 s. Separately, sampled precision accounting recorded six relay and eleven leaf map expiry/recovery episodes; sparse PPS intersections cannot erase those gaps.

This is a compelling fractional-second phase result for a causal two-link service with calibration frozen beforehand. It does not identify the physical cause of the empirical correction, calibrate receiver/cable delays, establish the UTC second association, or measure a delivered electrical output. Calibration portability across directions, boards, channels, temperature, and geometry is unknown. The native 62.5 ns counter grid does not prevent statistical sub-tick map coefficients, but neither does such estimation create a finer hardware output edge.

\subsection{Root and HOT capability, including regressions}
The 15 October checkpoint observed actual GM/HOT pairing in 24/24 sampled rows over 178.621 s. A separate monitor retained serving in 80/80 rows and HOT recovery in 64/76 standby rows \cite{legacy}. The same work retained three-GM instability and reduced the third unit to RC. No instrumented restart/takeover trace establishes failover latency or continuity. A later 179.717 s source-loss capture had 19 no-serving rows among 60 and no actual HOT. These observations demonstrate bounded synchronized standby and source withdrawal; current GM7 succession still needs physical loss/recovery and merge qualification.

\subsection{Heterogeneous relay and dynamic service}
The 14 October W12 release followed named direct RC1 paths in 108/117 and 106/115 frames, and authenticated depth-two RC2 paths in 91/116 frames in each of two treatments \cite{legacy}. Old-parent, rebind, SEARCH, HOLD and unavailable states remain. Publication required authenticated child demand and a usable upstream map; an observed expiry was consistent with the 360 s lease and 300 s renewal interval. Idle relays retained discovery/reception, with no measured energy saving. These functional estimated-TX results remain separate from the later W12 PPS population.

\subsection{ESP32-S3 event capture and heterogeneous phase delivery}
ESP32-S3 work separates counter resolution from physical event accuracy \cite{s3fine,s3low}. An extended 80 MHz capture counter supplies a 12.5 ns grid, while a bounded cross-clock bridge connects capture to the delivered coordinate. Narrow capture code, immutable model activation and pre-event evaluation avoid training an observation with its own or future correction. The RX-fall convention is empirical for a declared packet/PHY class; the delivered coordinate rounds to 50 ns. It is not a calibrated RF ADDRESS/END timestamp. TX remains a programmed nominal estimate.

The 22 October N30 run completed its fixed 660 s window with 657 eligible and three invalid PPS records, zero nominal missing, mean $-106.811\ns$, RMS $166.590\ns$ and maximum $1591\ns$ \cite{s3fine}. Four excursions above $1\us$ contribute 39.302\% of squared error. They remain scored; a good internal display residual cannot erase them. Reusing this population to select a later 51016 ns total offset makes it training for that offset, not its held-out validation. A fresh second-board transfer retained 48 eligible of 64 records, mean $+115.167\ns$, RMS $150.954\ns$ and maximum 362 ns. The sub-100 ns target and mean-zero calibration portability were not demonstrated.

The distinct 24 October study used W12 0.2.10 and T096 0.15.169, a declared 120 s warmup and two 300 s windows \cite{w12utc}. W12 followed BLE; its GNSS PPS was a witness and did not train that clock. The $+7671\ns$ W12 RX, $+4026\ns$ relay and $+10543\ns$ downstream RX adjustments were frozen before measurement. Of 600 logged W12 records, 522 were eligible: mean $+40.695\ns$, RMS $60.328\ns$, absolute P99/maximum 156/180 ns. All 78 source-unavailable records remain. The stable second window alone had 300/300 eligible and $38.312\ns$ RMS; this subset does not replace the full denominator.

A downstream T096 receiver yielded 451/600 eligible PPS comparisons, mean $-5.545\ns$, RMS $61.922\ns$ and P99/maximum 154/211 ns. Its 149 ineligible records comprise 109 no-map, 36 expired and four reference-ineligible observations. Separately, 58 of 120 selected emitted-timestamp comparisons at downstream RX measured $108.306\ns$ RMS and 226 ns maximum; 61 lacked follow-up matches and one lacked a PPS bracket. These observations test a heterogeneous timing path, while leaving TX-event uncertainty and output-pin accuracy unknown.

The source repair applies an integer-second translation, $T_{\mathrm{UTC}}=T_{\mathrm{rel}}+K\,10^9\ns$, preserving fractional phase. All 522 eligible W12 rows have identical relative and UTC phases; this fixes an earlier fractional-basis contamination, without independently establishing UTC association. The studies also do not identify multipath as the cause of timing tails. The later tail analysis distinguishes event-association, scheduling/aging, channel-dependent delay and oscillator hypotheses and proposes controlled interventions \cite{tails}. Calibration transfer needs fresh units, channels and placements, rather than retrospective centering of validation residuals.

\subsection{Extension into macOS}
The USB bridge associates device counter captures with host frame/Mach observations using causal markers and immutable source maps. This avoids treating USB request arrival time as the physical timing event. The fine application-clock path composes the USB map with the retained source coordinate and feeds a private \mbox{\texttt{chronyd -x}}; the operating-system clock is not steered. The chrony documentation describes the no-system-clock-control option \cite{chrony}; the project result concerns the private application path actually exercised.

The historical 15 October run measured $328.413\ns$ pre-ingest inverse forecast RMS on 9437 tuples and $414.351\ns$ settled shared-source agreement RMS on 182 valid comparisons, with 54 withheld and a 1315 ns maximum. Its full startup transient reached 32.072201 ms \cite{legacy,mac}. These remain separate from the newer comparison.

On 17 October, separate BLE-derived RC and local-GNSS Direct (DC) chains were observed concurrently for a fixed 180 s after-acquisition window on a shared USB bus, each requesting 32 samples/s \cite{mac}. DC is a separate Bluetooth-disabled GNSS/USB image. Each chain returned 5760 valid timing records; RC retained 5054 causal predictions, 694 warmup observations and 12 hardware rejects; DC retained 5370, 384 and six respectively. Bias-inclusive forecast RMS was $190.050\ns$ for RC and $267.448\ns$ for DC. These predict corrected source phase against prior SOF/Mach fits; they are not end-to-end UTC errors.

At 177 status points per chain, the user clock was delivered at 147 RC and 119 DC points, with 30 and 58 withheld. Separately admitted same-instant comparisons numbered 145 and 119, with agreement RMS $325.545/344.546\ns$ and maxima 909/831 ns. Cached chrony estimator readouts are a third population. All three metrics and sampled availability must remain distinct. Separate physical sources and an uncontrolled shared bus preclude attributing differences to BLE versus direct transport. Clean-bus, complete long-run and independent UTC validation remain open. A separate failed first-minute RC acquisition had no valid UTC replies; selecting the settled window does not remove that startup failure.

The result shows that useful relative timing can survive an inexpensive receiver-to-computer interface. It does not calibrate initial UTC offset or host/PHY bias, qualify the macOS system clock, or promise deterministic application scheduling. A timestamp query and a physical output action on a general-purpose OS have different latency and uncertainty budgets.

\section{Infrastructure topologies and extensibility}
\subsection{Managed cells and shallow relay trees}
A practical initial deployment is a managed timing cell: one independent root, one qualified standby where justified, a small set of powered RCs, and many passive clients. A star minimizes accumulated link uncertainty. A shallow relay tree extends around coverage obstacles while retaining a simple selected parent and diagnosable provenance. Powered infrastructure and battery endpoints can have different duty policies.

A corridor chain increases age, bias exposure, correlated failure and scheduling constraints. With an existing qualified backbone, several shallow cells may be preferable. Common upstream provenance permits analysis of their relationship, without proving phase alignment across bridges.

Multiple candidate parents can support mobility and local repair, provided switching fences the old path, reacquires the new one, and preserves the original source permission. RSSI and stream availability can guide attachment but are not timing-quality measurements. Parent dwell and hysteresis can reduce churn; a newly selected path may require different calibration. Fast packet connectivity must not masquerade as continuous clock service.

\subsection{Profile limits are separate from architectural reach}
Historical V14 permits eight route entries, three distinct RC positions, two consecutive RC intermediaries and four GM-control participants \cite{v14}. These are specification limits, not demonstrated deployment sizes. Current GM7 broadcast succession and RC7 dynamic path selection are successor designs \cite{current}. Root GR3 and relay RC7 retain 190-byte records; compatible coarse profiles remain required. Neither older topology limits nor later array capacities establish usable precision depth.

We define $H$ as received BLE links from root to endpoint: a direct leaf has $H=1$, one RC plus a leaf $H=2$, and two intermediaries plus a leaf $H=3$. The Nordic and later W12/T096 studies measure distinct $H=2$ paths, with different calibration and availability. No experiment here qualifies deep autonomous precision service. Larger-$H$ values below are sensitivity projections, not permission to exceed an implemented profile or controller limit.

\subsection{A transport-neutral clock interface}
Extensibility is most useful at the boundary between a clock contract and a transport adapter. The contract should return time, source scope, local and inherited age, uncertainty status, validity deadline, and discontinuity information. An adapter supplies actual physical-event evidence and its capture semantics. USB, Ethernet ingress, alternate MCUs, and future output hardware can then be evaluated against the same obligations without claiming identical performance.

An optional tunable oscillator or synthesizer could move from software clock prediction to physical frequency/phase control. It requires measured actuator gain, delay, quantization, saturation and a safe loop. Software-map convergence does not demonstrate quartz discipline, carrier coherence or calibrated PPS output. LoRa timing fallback and diagnostics are additional target capabilities, not evidence of BLE-equivalent timing. The current timezone extension similarly carries authenticated display metadata with its own presence lease; UTC mappings, precision payloads and source permission remain unchanged \cite{current}.

\section{Conditional hop-count and scalability analysis}\label{sec:scaling}
\subsection{Random error, systematic bias, and common source error}
Let the endpoint error after $H$ links be
\begin{equation}
e_H=e_0+e_{\mathrm{out}}+\sum_{i=1}^{H}(b_i+\epsilon_i)+e_{\mathrm{age}},
\end{equation}
where $e_0$ is source error, $e_{\mathrm{out}}$ is the final interface/output error, $b_i$ is residual link bias, $\epsilon_i$ is zero-mean link variation, and $e_{\mathrm{age}}$ is holdover error. With total mean $\mu_H$, the general RMS includes covariance:
\begin{equation}
\operatorname{RMS}_H^2=\mu_H^2+\operatorname{Var}\!\left(e_0+e_{\mathrm{out}}+e_{\mathrm{age}}+\sum_i\epsilon_i\right).
\end{equation}
Only under uncorrelated terms does the variance reduce to a sum. Shared source errors, calibration offsets, temperature, channel effects, and estimation histories can be correlated. Differential timing between two leaves may cancel common source error, but it does not cancel their distinct branches or output delays.

To anchor sensitivity to the measured scale without pretending to identify per-link noise, define the \emph{hypothetical effective increment}
\begin{equation}
s_{\mathrm{eff}}=62.503871/\sqrt{2}=44.1969\ns.
\end{equation}
This allocation treats the entire observed two-link RMS as two independent equal terms, with zero source, output, bias, and age contributions. Those assumptions are unverified. It is neither a measured per-link parameter nor a conservative confidence bound. Two illustrative extensions are $E_{\mathrm{ind}}(H)=s_{\mathrm{eff}}\sqrt H$ and $E_{\mathrm{corr}}(H)=s_{\mathrm{eff}}H$, the latter representing perfectly aligned equal random contributions at that assumed scale. It is a sensitivity alternative, not a second fit to the two-link result.
\begin{table}[htbp]\centering\small
\caption{Hypothetical RMS scaling in ns. Only the $H=2$ independent allocation is anchored to the measured RMS by construction; other entries are projections.}\label{tab:hops}
\begin{tabular}{@{}rrrr@{}}\toprule
BLE links $H$ & RC intermediaries & Independent model & Fully correlated model\\\midrule
1&0&44.2&44.2\\
2&1&62.5&88.4\\
3&2&76.6&132.6\\
4&3&88.4&176.8\\
8&7&125.0&353.6\\
16&15&176.8&707.2\\\bottomrule
\end{tabular}
\end{table}

Bias can dominate this optimistic random-error allocation. Adding a hypothetical same-sign 100 ns residual per link gives $\sqrt{H s_{\mathrm{eff}}^2+(100H)^2}$ ns: four links yield 409.7 ns and five 509.7 ns. Neither 100 ns/link nor a 500 ns illustrative budget is a measured deployment bound. Inverting the random-only curve into a large allowable hop count would ignore the unknown components and confuse RMS with a worst-case guarantee.

Unknown calibration/output terms cannot be zeroed in a deployment budget. The two-link experiment cannot separate source, reference, capture and relay components or identify the correlation model. Maximum usable depth requires deeper independent measurements.

\subsection{Source age and scheduling}
A relay's local update resets only its local estimator age. A simple age accounting model is
\begin{equation}
A_H=A_{\mathrm{ingress}}+\sum_{i=1}^{H}d_i,
\end{equation}
with each $d_i$ including the relevant capture-to-correction availability, mapping, and forwarding wait. These values require measurement. At a 16 Hz timing cadence the period is 62.5 ms. If, purely illustratively, each link consumed a full period, eight links would consume 500 ms before ingress age, missed events, or margins. This is not the measured per-hop delay, and the present local 500 ms rule is not an end-to-end age guarantee. Delayed corrections can span more than one event; off-phase scheduling can reduce some waits.

Slot capacity must satisfy a collision-domain constraint such as
\begin{equation}
K(t_{\mathrm{event}}+t_{\mathrm{guard}})+t_{\mathrm{control}}+t_{\mathrm{reserve}}\leq T_f.
\end{equation}
Here $K$ is active transmitting trains in the domain, not passive receivers. Controller activity limits, upstream/downstream receive conflicts, coexistence, and guards may bind before packet duration does. Spatial reuse requires evidence about interference, hidden terminals, and channel scheduling; multiplying a single-cell capacity by area is unjustified without it.

\subsection{Airtime, fan-out, and control load}
A periodic broadcast can serve many listeners without one precision exchange per listener. For a 190-byte record at 16 Hz, record bits alone require $190\times8\times16=24320$ bit/s, or 2.432\% of a 1 Mbit/s raw link rate. This lower-bound serialization excludes PHY/link overhead, discovery, auxiliary scheduling, redundancy, control and guards. Ten trains sharing one radio resource require at least 24.32\% on this simplistic basis; that does not demonstrate a feasible schedule or collision performance. Historical 212-byte and other formats need separate accounting; 212 bytes is not the current selected precision record.

Broadcast TX work is approximately independent of passive client count; admission state and renewal work are not. With $N_a$ admitted children and interval $T_r$, mean request rate is $N_a/T_r$: a hypothetical 1000 children renewing every 300 s produce 3.33 requests/s before retries/bursts. This is not demonstrated capacity. Memory, admission serialization, startup and source-change bursts can bind before precision airtime.

\subsection{Availability scales differently from RMS}
If each link were independently usable with probability $p$ at a defined observation instant and every link were required simultaneously, path usability would be $p^H$. For $p=0.99$, ten and twenty links give 90.4\% and 81.8\%; for $p=0.999$, twenty give 98.0\%. These are illustrative probabilities, not estimates from the selected experiments. Holdover changes the event definition and correlated outages invalidate independence. The measured 72.425\% relay scoreability includes reference validity and cannot be substituted for a per-link packet-success probability.

Redundant parent choices may improve coverage but incur acquisition and continuity costs; two correlated paths are not independent redundancy. Source power loss, interference, a shared GNSS obstruction, or a common software defect can affect all routes. A viable network must report both the conditional error distribution when usable and the time/observation denominator over which it is usable. The effective deployment depth is limited jointly by timing, age, schedule, availability, trust, and supported profile capacity.

\section{Power sampling, accuracy, and economics}
\subsection{Why the periodic stream is well positioned}
Periodic timing offers predictable reception opportunities, repeated frequency observations, and a way to amortize acquisition and clock fitting. Multiple clients can share a stream, and a receiver can potentially accumulate sufficient observations in a scheduled window before sleeping. Complete corrections avoid relying on incomplete fragment assembly for each useful update. Demand-driven relay publication can avoid serving an empty branch.

These properties motivate energy-aware synchronization, but the experiments do not establish an optimum rate or energy per valid timestamp. The 19 October awake T096 pilot selected 125 ms training windows around a nominal 16 Hz grid for six separate runs with 5/10/20 s HOLD intervals \cite{power}. Of 120 main windows, 42 refreshed and 78 were empty. RF and the precision counter stayed on: the 2.439/1.235/0.621\% scheduled training fractions are neither radio duty nor energy savings. HOLD PPS RMS ranged from 80.259 to 2173.739 ns across runs; only one replaced every acquisition-fit pair. Missing BLE bins, failed priming and all PPS eligibility states remain. No causal ranking of schedules follows from these different histories.

\subsection{Measured frozen-clock holdover}
The 30 October W12 study tests BLE-trained frozen clock copies against later accepted BLE witnesses \cite{holdover}. At entry, each copy freezes phase, selected rate and residual slew; later BLE scores it without correcting it. For run $j$ and tested age $a$, define the prefix maximum $M_j$ and, for a fully witnessed $n$-run population, its quantile:
\begin{equation}
 M_j(a)=\max_{i:\,s_j<\ell_i\leq\ell_{k_j(a)}}|\widehat T_j(\ell_i)-R_i|,
 \qquad Q_{.99}(a)=M_{(\lceil .99 n\rceil)}(a),
\end{equation}
where $k_j(a)$ is the first actual witness at or after age $a$. Its overshoot and all earlier errors remain; a bad prefix cannot recover through a favorable endpoint. Each age retains up to 128 outcomes and requires at least 100 scored runs. Passing permits at most 1\% failed or unknown outcomes. Unknown expiry counts against passing and cannot supply the scored minimum. Duration requires every preceding tested age to pass the strict $1000\ns$ threshold.

The best fully replayed C-1 rolling window supports 64 s on 128 prefixes, P99/maximum 760/879 ns. A separate fixed mature long-run cohort supports 75 s, P99 983 ns, and fails at 76 s, P99 1028 ns; it is an older cohort, not refinement of the same rolling window. C-2's first normal capture supports at least 48 s, P99/maximum 606/880 ns. Later capture-end rolling durations are 24 s on both devices. Historical longer holds include failures, including a descriptive 51-run 96 s cohort with 4110 ns P99/maximum, below the 100-run readiness minimum. Overlapping starts share history; the empirical P99 supplies no independent-trial confidence or continuous gap-interior bound.

Production HOLD uses a causal, channel-aware BLE rate selector with a capped 5 ppb entry adjustment and fallback; temperature steering remains disabled. An earlier thermal trial reached 3663 ns maximum versus 362 ns for its paired baseline. The separate shadow study preserves that regression. These measurements support a useful sampled holdover capability within particular device/time windows. They do not grant downstream source authority, establish a physical GM-transition result, or measure sleep power.

\subsection{An explicit holdover and energy model}
For a proposed sleep interval $T$, a deterministic budget could be
\begin{equation}
E(T)\leq E_{\mathrm{start}}+B_yT+\tfrac12 B_{\dot y}T^2+E_{\mathrm{readout}},
\end{equation}
where $B_y$ bounds residual fractional-frequency error and $B_{\dot y}$ bounds its change. These must be experimentally bounded, including temperature and aging. With a hypothetical remaining allocation of 500 ns and no drift term, 1 ppm permits 0.5 s, 0.1 ppm permits 5 s, and 10 ppb permits 50 s. None is an observed safe sleep interval or authority extension. Source permission may expire before an oscillator budget is exhausted.

A first-order power model is
\begin{equation}
\overline P=D P_{\mathrm{rx}}+(1-D)P_{\mathrm{sleep}}+f_{\mathrm{wake}}E_{\mathrm{acq}}+P_{\mathrm{compute}}+P_{\mathrm{other}}.
\end{equation}
At short windows, repeated synchronization acquisition can erase radio-off savings. Fewer samples can reduce estimator quality; excessive filtering can hide temperature-induced rate changes. Conversely, always listening need not improve accuracy if dominant bias is unchanged. The relevant optimization is energy per application-usable event subject to error, availability, and source-age constraints, rather than duty cycle alone.

The next study should sweep reception window, update rate, source/relay role, temperature, RF load, and motion while independently measuring current and delivered error throughout the entire blackout. It must retain failed acquisitions, wake latency, all unavailable intervals, and observer effects. A prospective held-output trace is required; retrospective fitting across future samples cannot qualify a quiet interval. This is especially important for the project's deferred $\pm500\ns$ total held-time objective at at most 50\% listening.

\subsection{Economics without unsupported prices}
The economic hypothesis is that a small number of infrastructure clocks can distribute useful time to existing BLE-capable endpoints, reducing endpoint GNSS hardware, antennas, cabling, or specialized network interfaces. It is not yet a measured cost advantage. A complete model is
\begin{equation}
C_{\mathrm{total}}=C_{\mathrm{sources}}+C_{\mathrm{relays}}+N C_{\mathrm{endpoint}}+C_{\mathrm{install}}+C_{\mathrm{energy}}+C_{\mathrm{maint}}+C_{\mathrm{qual}}.
\end{equation}
Precision integration includes controller access, oscillator quality, calibration, provisioning and maintenance. Coarse event-correlation clients and calibrated precision endpoints may occupy different cost tiers. Compare cost per qualified usable endpoint with the named alternative; no universal savings percentage is measured.

\section{A future PTP-to-PToB edge bridge}\label{sec:ptp}
The most direct extension to existing infrastructure is a PTP-disciplined ingress clock that publishes PToB time into a local BLE cell. A hardware-timestamped Ethernet interface would discipline a local clock, with a calibrated mapping to the radio counter. The root would retain upstream source identity and validity. The later robotics study develops this into a proposed PHC-to-radio bridge, with domain, ancestry, event and quality fields \cite{robotics}. Its first target is wireless acquisition timestamps comparable with wired events. No PTP/gPTP bridge, transparent carriage, TSN traffic schedule or physical actuation is demonstrated by the selected evidence.

The bridge must preserve domain, timescale, epoch, UTC/leap handling, traceability, source changes and loss behavior. Its Ethernet-to-radio cross-timestamp needs measured latency and uncertainty; a callback's wall-clock read is insufficient. Upstream changes must fence incompatible calibration and authority. Two bridges sharing one PTP source do not provide independent source redundancy.

For a deterministic end-to-end bound, a proposed allocation is
\begin{equation}
B_{\mathrm{edge}}\leq B_{\mathrm{PTP}}+B_{\mathrm{cross}}+\sum_i B_{\mathrm{BLE},i}+B_{\mathrm{hold}}+B_{\mathrm{output}}.
\end{equation}
Every component requires a supported bound. Root-sum-square combination is appropriate only for explicitly modeled random components and does not remove bias. If any required bound is unknown, the total remains unknown. A calibrated PTP backbone cannot by itself calibrate the wireless edge.

The first gateway should follow one upstream authority and demote on its loss. Reverse PTP export requires separate conformance and loop prevention: mesh-derived time cannot re-enter upstream as an independent superior source. IEEE edition, profile, delay mechanism and time properties belong to that specific interoperability design.

Validation proceeds from cross-counter calibration through direct BLE, one relay, source switching and holdover. Independent simultaneous ingress/output measurements must separate common-mode and differential error. SMPTE profiles and physical media-clock/genlock outputs need their own validation.

\section{Application viability}
\begin{table}[htbp]\centering\small
\caption{Application-oriented assessment. Every precision requirement is application-specific.}\label{tab:apps}
\begin{tabular}{@{}L{0.22\linewidth}L{0.35\linewidth}L{0.37\linewidth}@{}}\toprule
Application & Plausible value & Remaining decisive evidence\\\midrule
Monitoring / sensor fusion & Source-aware machine, environmental and robot acquisition timestamps & Capture-delay calibration, coexistence, mobility, declared outage budget and energy per usable event\\\addlinespace
Industrial control & Prospective scheduled sensing or actuation & Worst-case clock/latency bounds, physical output, continuity and safe failure response\\\addlinespace
Broadcast & Accessory telemetry, production metadata; prospective timecode adapters & Bridge/profile semantics, exact endpoint interoperability, media output and redundancy\\\addlinespace
Positioning & Prospective surveyed beacon and receiver-clock solve & Differential anchor calibration, detector/path bias, geometry, motion and independent spatial truth\\\bottomrule
\end{tabular}
\end{table}

Event correlation is the strongest initial use case because an application can explicitly reject unavailable or uncertain timestamps and retain source provenance. Higher-integrity systems should propagate quality with the event data, rather than converting all samples to an apparently uniform UTC field. Two events with different roots or source epochs require a qualified relationship before their timestamps can be compared precisely.

The broadcast opportunity is complementary to a facility backbone: extending time to accessories and sensors. The timecode study separates a common source coordinate from exact frame/sample mapping and device-specific adapters \cite{timecode}. Bluetooth control of a generator, timecode metadata reception, LTC output and genlock are different interfaces. Their existence in commercial products does not establish PToB interoperability. A relative software clock suitable for logs is not a qualified media sample clock.

The indoor PNT study likewise proposes using event time beneath heterogeneous sensor fusion \cite{pnt}. A shared clock can reduce acquisition-time ambiguity without creating a positioning observation. The distinct PToB-IPS proposal adds surveyed beacons, aligned multi-rate advertising slots and a joint position/receiver-clock solve \cite{ips}. One-way propagation changes by about 0.30 m/ns: a hypothetical uncompensated 60 ns timing bias corresponds to about 18 m of range bias before geometry, rather than centimetre positioning. Surveyed delay, differential beacon timing, RF detector bias and motion need separate constraints; averaging cannot identify a constant unknown path delay. Neither paper reports an implemented positioning service.

A heterogeneous service can offer different timing tiers, provided capability discovery preserves quality: recognizing the same root does not make a coarse client calibrated.

\section{Limitations and a validation program}
The evidence comes from a small set of laboratory and field-oriented prototypes, finite treatments, different software generations and incompletely characterized environments. Some results are implementation-level functional checks; independent source-to-output qualification remains incomplete. There is no large-scale deployment, lifetime energy study, controlled mobility campaign, adversarial security validation or calibrated complete source-to-electrical-output budget. One-way empirical compensation leaves propagation and detector/reference effects entangled. At roughly 3.3 ns per metre of free-space path, changing geometry can matter at a tens-of-nanoseconds objective; this scale is not a measurement of the tested path.

The present synthesis does not show that the best Nordic relay, HOT, W12, holdover and macOS results coexist in one configuration. Small correlated cohorts cannot establish rare-event reliability; sample P99 is not a population confidence bound. Calibration learned on one packet class or unit may fail on another. Firmware regressions, source loss, gaps and incomplete captures remain. Hashes bind identity, not measurement truth. Public selected data support numerical replay, while excluded raw acquisition and calibration prevent complete independent reconstruction.

A staged validation program follows directly from the viability model:
\begin{enumerate}
\item \textbf{Metrology:} freeze calibration; measure named output/capture points against an independent calibrated reference, including cable, detector and counter uncertainty.
\item \textbf{Composition:} test root/HOT, heterogeneous relays and host delivery together. Withdraw authority while phase remains predictable; change epochs and reorder follow-ups. Count false-valid output separately from phase error and unavailable service.
\item \textbf{Control:} inject primary loss, asymmetric reception, arrivals, restart, partition and merge. Score dual serving, HOT tenure, convergence, discontinuity and lost service.
\item \textbf{Depth/density:} test direct, one- and two-intermediary paths. Vary transmitters independently of listeners; measure occupancy, controller limits and admission bursts.
\item \textbf{Energy/environment:} pair current with continuously observed held-output error through quiet windows, temperature, coexistence and mobility.
\item \textbf{PTP/application:} validate ingress clock mapping, properties and failover, then the named application endpoint against its profile and instrumentation.
\end{enumerate}
These integrated qualification steps remain proposed. The awake pilot and sampled holdover study supply partial evidence for step 5, not measured sleep energy or a continuous blackout bound. Future campaigns should preregister error, availability and recovery thresholds, freeze calibration and images, then score all held-out opportunities. A smaller conditional RMS with more unavailable events is not automatically an improvement.

\section{Conclusion}
PToB has progressed from isolated clock fitting to measured causal timing across Nordic relays, ESP32-S3 receivers and a heterogeneous relay path, with separate host-delivery and sampled holdover evidence. Relative PPS RMS near 60 ns is achievable in selected populations; 1591 ns tails, source-unavailable observations and shorter later hold windows are equally material results. Bounded HOT observations and current GM7 mechanisms support further source-continuity research, without establishing physical failover.

The evidence supports shallow managed timing cells more strongly than deep autonomous meshes. Broadcast fan-out is favorable, but depth couples bias, age, availability and controller scheduling. The awake pilot and frozen-clock study make energy-aware operation testable, while measured sleep savings and safe low-duty precision remain open. A hardware-associated PTP ingress could reuse the source/clock-map contract after its calibration and semantics are validated.

The resulting research claim is specific: heterogeneous BLE time sharing is credible at several measured interfaces, especially for source-aware event correlation. Industrial and broadcast qualification depends on composed continuity, independent output accuracy, dynamic topology and energy evidence.

\section*{Author contributions and AI-assistance disclosure}
Jonathan Yantis supplied the protocol, hardware, experimental direction, engineering decisions and research objectives. OpenAI Codex assisted implementation, tooling, evidence analysis, literature checking, manuscript synthesis, analytical models, LaTeX preparation and document verification. The human author retains responsibility for all claims and any submission. No AI system is an author.

\section*{Data and source availability}
The standalone LaTeX embeds its bibliography and vector architecture figure. Source, selected data, dated provenance and standard-library offline verifiers are maintained in the \href{https://github.com/yantisj2/PToB}{PToB repository}; Appendix~\ref{app:reproduce} maps the evidence. The 30 October edition manifest binds this source/PDF pair and its inputs; earlier editions remain in Git. Papers, figures and original selected/derived data use CC BY 4.0, credit Jonathan Yantis; original executable tools and ordinary documentation use MIT. Third-party terms remain. Private enrollment, images and raw operational captures are excluded. A durable archival identifier has not been assigned; repository availability is not an independent calibration or archival deposit.

\appendix
\section{Evidence map and reproducibility}\label{app:reproduce}
\begin{table}[htbp]\centering\small
\caption{Public repository anchors. Older private archive locators are preserved only as provenance, not required reproduction inputs.}
\begin{tabular}{@{}L{0.24\linewidth}L{0.70\linewidth}@{}}\toprule
Claim family & Repository collection / verifier\\\midrule
Historical PLL and relay & \path{evidence/papers/oscillator-pll-2026-10-03/}; \path{evidence/v14/2026-10-12/}; \path{tools/verify_evidence.py}; \path{tools/verify_v14.py}\\
Original synthesis / models & \path{evidence/research/2026-10-16/}; \path{tools/verify_research.py}\\
USB RC / DC & \path{evidence/usb/2026-10-17/}; \path{tools/verify_usb_results.py}\\
Awake pilot & \path{evidence/power/2026-10-19/}; \path{tools/verify_power.py}\\
N30 fine / W12 relay & \path{evidence/s3/2026-10-22-fine/}; \path{evidence/s3/2026-10-24-w12-utc-basis/}; \path{tools/verify_s3_fine_paper.py}; \path{tools/verify_w12_utc_results.py}\\
Frozen-clock holdover & \path{evidence/w12/2026-10-30-holdover/}; \path{papers/mcxo-holdover-duration/tools/verify.py}\\
This edition & \path{papers/editions/viability-2026-10-30.json}; \path{tools/verify_viability_paper.py}\\\bottomrule
\end{tabular}
\end{table}

From the repository root, \texttt{python3 tools/verify\_viability\_paper.py} checks edition identities, claims and public anchors; \texttt{--replay} reruns the supporting numerical verifiers. Python 3.9+ and recorded public Git history suffice, without hardware or credentials. Public Tectonic 0.17.0 compiles the standalone source with ordinary LaTeX packages, no external figures/bibliography and no shell escape. arXiv processing and author metadata review remain submission-stage checks.

The hypothetical hop, availability and holdover tables regenerate from the retained \path{evidence/research/2026-10-16/viability/derive.py}. They are unchanged sensitivity calculations, not fitted hardware bounds. Literature was rechecked on 30 October 2026 against primary papers, publisher metadata and official documentation. Where only publisher abstracts were available, no detailed experimental reconstruction is claimed. Standards editions are explicit context references, not assertions of latest-edition conformance.

\begin{thebibliography}{99}\small\raggedright
\bibitem{ftsp} M. Mar\'oti, B. Kusy, G. Simon, and \'A. L\'edeczi, ``The Flooding Time Synchronization Protocol,'' SenSys, 2004, pp. 39--49. \href{https://doi.org/10.1145/1031495.1031501}{doi:10.1145/1031495.1031501}.
\bibitem{cheep} S. Sridhar, P. Misra, G. S. Gill, and J. Warrior, ``CheepSync: A Time Synchronization Service for Resource Constrained Bluetooth Low Energy Advertisers,'' IEEE Communications Magazine, vol. 54, no. 1, pp. 136--143, 2016. \url{https://arxiv.org/abs/1501.06479}.
\bibitem{bluesync} F. Asgarian and K. Najafi, ``BlueSync: Time Synchronization in Bluetooth Low Energy With Energy-Efficient Calculations,'' IEEE Internet of Things Journal, vol. 9, no. 11, pp. 8633--8645, 2022. \href{https://doi.org/10.1109/JIOT.2021.3116921}{doi:10.1109/JIOT.2021.3116921}.
\bibitem{micro} R. Ohara et al., ``MicroSync: Sub-Micro Second Accuracy Wireless Time Synchronization Service,'' IEEE Access, vol. 12, pp. 124478--124494, 2024. \href{https://doi.org/10.1109/ACCESS.2024.3446668}{doi:10.1109/ACCESS.2024.3446668}. \url{https://hdl.handle.net/20.500.14094/0100491658}.
\bibitem{paenergy} S. Gautam, S. Kumar, and P. P. Singh, ``A Comparative Analysis of Periodic Advertisements and Connections in Bluetooth Low Energy,'' IEEE Wireless Communications Letters, vol. 15, pp. 96--100, 2026; online 14 October 2025. \href{https://doi.org/10.1109/LWC.2025.3621352}{doi:10.1109/LWC.2025.3621352}. Publisher abstract consulted.
\bibitem{parelay} S. Gautam, S. Kumar, and E. E. Tsiropoulou, ``Unlocking Ultralow-Power Bluetooth Low Energy Relays With Periodic Advertisements,'' IEEE Internet of Things Journal, vol. 13, no. 12, pp. 27717--27735, 2026. \href{https://doi.org/10.1109/JIOT.2026.3680032}{doi:10.1109/JIOT.2026.3680032}. \href{https://github.com/sukgautamgit/DPCS-SimultaneousTx-Rx-BLE-Demo}{Author relay artifact inspected}; full paper unavailable.
\bibitem{ptp} IEEE, \emph{IEEE 1588-2019: IEEE Standard for a Precision Clock Synchronization Protocol for Networked Measurement and Control Systems}, 2020. \url{https://standards.ieee.org/ieee/1588/6825/}.
\bibitem{gptpstandby} IEEE, \emph{802.1ASdm-2024: Hot Standby and Clock Drift Error Reduction}, published 4 October 2024, \url{https://1.ieee802.org/tsn/802-1asdm/}; \emph{802.1ASed-2026: Fault-Tolerant Timing with Time Integrity}, published 25 March 2026, \url{https://1.ieee802.org/tsn/802-1ased/}. Official scope records.
\bibitem{securetime} T. Mizrahi, \emph{Security Requirements of Time Protocols in Packet Switched Networks}, RFC 7384, October 2014, Secs. 5.1.1--5.1.2, 5.1.5. \url{https://www.rfc-editor.org/rfc/rfc7384.html}.
\bibitem{delaysecurity} L. Narula and T. E. Humphreys, ``Requirements for Secure Clock Synchronization,'' accepted manuscript, 2018. \url{https://arxiv.org/abs/1710.05798v3}. \href{https://doi.org/10.1109/JSTSP.2018.2835772}{doi:10.1109/JSTSP.2018.2835772}.
\bibitem{smpte} SMPTE, \emph{ST 2059-2:2021: SMPTE Profile for Use of IEEE-1588 Precision Time Protocol in Professional Broadcast Applications}, fixed-edition reference. \url{https://github.com/SMPTE/st2059-2}. Repository revision drafts are distinct from this edition.
\bibitem{ble} Bluetooth SIG, \emph{Bluetooth Core Specification 6.1}, Vol. 6, Part B, Link Layer Specification, fixed-edition reference. \href{https://www.bluetooth.com/wp-content/uploads/Files/Specification/HTML/Core-61/out/en/low-energy-controller/link-layer-specification.html}{Official Link Layer specification}.
\bibitem{rpl} T. Winter et al., ``RPL: IPv6 Routing Protocol for Low-Power and Lossy Networks,'' RFC 6550, 2012. \url{https://www.rfc-editor.org/rfc/rfc6550.html}.
\bibitem{chrony} chrony project, \emph{chronyd(8)}, version 4.9, option \texttt{-x}. \url{https://chrony-project.org/doc/4.9/chronyd.html}.
\bibitem{legacy} J. Yantis, \emph{PToB Viability: Rooted BLE Timing Meshes for Heterogeneous Edge Networks}, 16 October 2026 research draft. \repo{papers/ptob-viability-preprint.tex}{Preserved original source}; \repo{evidence/research/2026-10-16/provenance/ptob-viability-evidence-20261016.json}{claim/input provenance}. Includes retained 14--15 October HOT and heterogeneous-relay summaries; underlying private captures excluded.
\bibitem{osc} J. Yantis, \emph{Causal Software Clock Discipline over Wireless Physical Events}, 3 October 2026 research draft. \repo{papers/oscillator-pll-research-20261003.tex}{Manuscript}; \repo{evidence/papers/oscillator-pll-2026-10-03/README.md}{selected evidence}.
\bibitem{v14} J. Yantis, \emph{Passive Precision Time over BLE: V14 Boundary Clocks, Causal Timestamping, and Prospective One-Relay Validation}, 12 October 2026 research draft. \repo{papers/ptob-v14-preprint.tex}{Manuscript}; \repo{evidence/v14/2026-10-12/README.md}{one-relay data and reproduction}.
\bibitem{current} J. Yantis, \emph{Passive Precision Time over BLE: V14 Candidate Service and Profile Proposal}, 30 October 2026 author-review draft. \repo{papers/ptob-v14-sig-proposal.tex}{Current design}; \repo{docs/t096/gm7-implementation.md}{GM7 implementation and limits}. No new hardware qualification.
\bibitem{s3fine} J. Yantis, \emph{Precision Bluetooth Timing on ESP32-S3: Causal Clock Discipline and Hardware Event Capture}, 22 October 2026 research draft. \repo{evidence/s3/2026-10-22-fine/README.md}{Selected observations and reproduction}.
\bibitem{s3low} J. Yantis, \emph{Low-Level Timing on ESP32-S3: Event Capture, Clock Bridges, and Causal Validation}, 24 October 2026 research draft. \repo{papers/ptob-esp32s3-low-level-timing.tex}{Implementation study}.
\bibitem{w12utc} PToB project, \emph{W12 baseline after the T096 UTC-basis repair}, 24 October 2026. \repo{evidence/s3/2026-10-24-w12-utc-basis/README.md}{Selected PPS/TX rows, results and reproduction}.
\bibitem{tails} J. Yantis, \emph{Reducing Timing Tails and Multipath Sensitivity in an Operating BLE Precision Mesh}, 26 October 2026 research draft. \repo{papers/ble-tail-multipath/README.md}{Source and analytical supplement}. No measured multipath attribution.
\bibitem{mac} J. Yantis, \emph{Beyond USB Packet Arrival: Frame-Identified Time Transfer to macOS}, 17 October 2026 research draft. \repo{papers/usb-macos-time-transfer-preprint.tex}{Manuscript}; \repo{evidence/usb/2026-10-17/README.md}{RC/DC data}. Earlier 15 October CSVs remain in \repo{evidence/research/2026-10-16/README.md}{the historical collection}.
\bibitem{power} J. Yantis, \emph{How Much Timing Information Is Enough? Awake Hardware Evidence and Wake/Listen Energy Models for PToB}, 19 October 2026 research draft. \repo{evidence/power/2026-10-19/README.md}{Selected observations and models}.
\bibitem{holdover} J. Yantis, \emph{Empirical Sub-Microsecond Holdover on BLE-Trained Software Clocks}, 30 October 2026 research draft. \repo{papers/mcxo-holdover-duration/README.md}{Manuscript and replay}; \repo{evidence/w12/2026-10-30-holdover/README.md}{selected archive and provenance}.
\bibitem{robotics} J. Yantis, \emph{A Shared Time Base for Robotic gPTP and TSN Systems}, 22 October 2026 research draft. \repo{papers/ptob-gptp-tsn-robotics-preprint.tex}{Proposed edge adaptation and validation programme}.
\bibitem{timecode} J. Yantis, \emph{A Shared Time Base for Bluetooth Timecode Ecosystems}, 22 October 2026 research draft. \repo{papers/ptob-bluetooth-timecode-preprint.tex}{Proposed compatibility layers}.
\bibitem{pnt} J. Yantis, \emph{A Shared Time Base for Indoor PNT}, 26 October 2026 research draft. \repo{papers/indoor-pnt/README.md}{Source and analytical supplement}.
\bibitem{ips} J. Yantis, \emph{PToB-IPS: Toward GNSS-like Indoor Positioning with Surveyed BLE Timing Beacons}, 26 October 2026 theoretical research draft. \repo{papers/indoor-ips/README.md}{Source, assumptions and reproduction}.
\end{thebibliography}
\end{document}
